\section{Serving and training configuration}
\label{app:config}

\subsection{Per-family serving flags}

These are the configurations under test, including documented serving flags and the
adapter/schema interventions identified in our experiments. The rightmost column summarizes
observed limitations or failure symptoms; it is not a uniform flag-omission ablation.

{\small
\begin{longtable}[]{@{}p{.17\linewidth}p{.48\linewidth}p{\dimexpr.35\linewidth-4\tabcolsep\relax}@{}}
\caption{Serving configurations and observed symptoms (vLLM 0.27.1).}
\label{tab:serving}\\
\toprule
Family & Flags / intervention & Observed symptom \\
\midrule
\endfirsthead
\toprule
Family & Flags / intervention & Observed symptom \\
\midrule
\endhead
all & \texttt{-{}-enable-auto-tool-choice} & HTTP 400 if omitted \\
Qwen2.5 (Instruct) & \texttt{-{}-tool-call-parser hermes} & silent empty \texttt{tool\_calls} \\
Qwen2.5-Coder & \emph{no working stock configuration}; we install the community
  \texttt{<tools>} parser plus its chat template & silent, at every scale \\
Llama 3.1 / 3.2 & \texttt{-{}-tool-call-parser llama3\_json} \newline
  \texttt{-{}-chat-template} \path{tool_chat_template_llama3.1_json.jinja} \newline
  plus \texttt{"strict": true} in the schema & malformed arguments; without
  \texttt{strict} the terse baseline has 23\% wrong-tool calls \\
Mistral (HF format) & \texttt{-{}-tokenizer-mode hf -{}-config-format hf -{}-load-format hf}
  \newline \texttt{-{}-tool-call-parser mistral} \newline
  \texttt{-{}-chat-template\allowbreak\ tool\_\allowbreak chat\_\allowbreak template\_\allowbreak mistral\_\allowbreak parallel.jinja} \newline
  \texttt{tool\_call\_id} must be exactly 9 characters & HTTP 400; the officially
  recommended parallel template \emph{raises} the error rate from 2\% to 42\% \\
DeepSeek-Coder & none available: the chat template never injects \texttt{tools} & n/a \\
\bottomrule
\end{longtable}
}

Missing auto-tool-choice support and the tested Mistral configurations produce explicit
request errors. Other mismatches can return HTTP 200 with empty \texttt{tool\_calls},
indistinguishable at that field from declining to call a tool. The released 98-line
\texttt{preflight\_toolcall.py} checks a canonical request for a non-empty call with the
right name and parseable arguments, then repeats under \texttt{tool\_choice: required}
as a positive control. This is a screening check, not a guarantee of compatibility on all tasks.

\subsection{Training configuration}

The historical 1.5B runs use the following configuration; they motivate the investigation but are
not the formal parser-repair control:

{\small
\begin{longtable}[]{@{}p{.27\linewidth}p{\dimexpr.73\linewidth-2\tabcolsep\relax}@{}}
\caption{Training configuration. All arms share every value below; arms differ only in the
protocol and, for the historical runs, the algorithm and seed.}\label{tab:trainconfig}\\
\toprule
Setting & Value \\
\midrule
\endfirsthead
\toprule
Setting & Value \\
\midrule
\endhead
Model & Qwen2.5-Coder-1.5B-Instruct \\
Framework & verl 0.9.0 (vLLM 0.27.1 rollout backend) \\
Algorithm & GRPO ($\times$2 seeds) and RLOO ($\times$1 seed) \\
Trajectories per step & \texttt{train\_batch\_size} 16 $\times$ \texttt{rollout.n} 8 = 128 \\
Steps & 150 for the historical runs and formal 7B control; 10 / 3 / 75 for the historical
diagnostic arms in Appendix~\ref{app:extra} \\
Training data & KodCode, decontaminated pool of 7669 items \\
Held-out evaluation & EvalPlus: HumanEval+ 319, MBPP+ 677 \\
Denominators & 540 (multi channel), 454 (repair channel), after removing 2 known-defective items \\
Agent loop & \texttt{tool\_agent} (FC) or \texttt{react\_agent} (ReAct) \\
Tool & \texttt{run\_tests}, single \texttt{code} argument, sandboxed \texttt{pytest} \\
System prompt & \texttt{FC\_MANDATORY} / ReAct mandatory --- matched strength (Appendix~C) \\
Hardware & single A800 80\,GB \\
\bottomrule
\end{longtable}
}

The formal control of \S\ref{sec:rl} instead uses Qwen2.5-Coder-7B-Instruct, GRPO, LoRA rank and
alpha 32 over all linear modules, seed 0, 150 steps, and the same 16 prompts $\times$ 8 rollouts
per step in both arms. Both use the mandatory FC prompt and fixed verifier, retain complete rollout
text, and evaluate the same 542 paired multi-turn items at steps 0/30/60/90/120/150. Evaluation
uses a shared parser-independent programmatic-feedback protocol: the current LoRA is queried
through verl's native rollout RPC at temperature 0, for at most four turns and 1,024 tokens per
turn; generated code is executed against EvalPlus tests and the observation is appended as a user
message. Native RPC metadata must match the expected global weight step. The sole arm
difference is \texttt{multi\_turn.format=hermes} versus the registered
\texttt{qwen2\_5\_coder} parser. Full-parameter 7B RL did not fit the single 80\,GB card, so both
arms use the same parameter-efficient configuration; this preserves the arm contrast but prevents
a strict comparison to the historical full-parameter curves.

\subsection{Tool-call instrumentation}

verl 0.9.0's step metrics contain \textbf{no} tool-call count. The metric named
\texttt{timing\_s/agent\_loop/tool\_calls} is \emph{elapsed seconds}; reading it as a count is
a hard error and we flag it because the name invites exactly that.

We therefore instrumented the two agent-side call sites directly, each writing one labelled
line per call:

\begin{itemize}
\tightlist
\item \texttt{CodeTool.execute} --- the function-calling path
\item \texttt{ReActAgentLoop.\_run\_tests} --- the ReAct path
\end{itemize}

\textbf{The counter must not be placed in \texttt{sandbox.run\_tests}.} The reward function
evaluates the final program through that same function, so a counter there conflates reward
evaluation with agent tool use. We hit this and moved the counter; the sandbox's own counter
was renamed \texttt{sandbox\_exec\_count} so the two can never be confused again.

In the historical diagnostic run, the FC counter file has \textbf{0 lines}; the ReAct arm's has \textbf{788}, all
carrying the \texttt{react} label, so the provenance is uncontaminated. Three mutually
independent signals agree --- \texttt{num\_turns} pinned at its minimum of 2, tool time
exactly 0.00\,s, and a dedicated counter at 0.

\subsection{Constructing the repaired-FC training arm}

The control initially appeared blocked at three layers:

\begin{enumerate}
\tightlist
\item \texttt{multi\_turn.format=hermes} selects \textbf{verl's own} \texttt{ToolParser}
  registry, which is independent of vLLM's --- so installing a vLLM parser plugin has no
  effect on the training path.
\item A verl-side parser we wrote does recover calls, but \textbf{none of them are the tool}:
  36/52 target the task function, 16/52 carry executable code elsewhere in the body. Name
  normalisation cannot repair this, because no call carries code to normalise.
\item The 1.5B checkpoint emits no correctly named call for that parser to accept, and 7B
  full-parameter training does not fit a single 80\,GB card.
\end{enumerate}

The first two observations determine the intervention; the third determines the scale and LoRA
configuration. We registered the parser inside verl and ran the matched 7B pair. Thus the earlier
statement that no executable repaired-FC arm existed was a description of the configuration
surface, not a permanent limitation; \S\ref{sec:rl} reports the completed control.

\subsection{The path to the repaired-FC training arm}
\label{app:norepair}

We encountered four successive obstacles while constructing the control. Each attempt is a
measurement, so we retain the record; all were eventually resolved by registering the parser in
verl and moving to matched 7B LoRA arms.

\textbf{(i) A vLLM parser plugin does not reach the training loop.} verl's AgentLoop performs its
own tool-call parsing (\texttt{verl/experimental/agent\_loop/tool\_parser.py}, registered as
\texttt{hermes}, adapted from vLLM v0.9.1). \texttt{multi\_turn.format=hermes} selects \emph{that}
parser, not the server's, so attaching a plugin to the vLLM rollout server changes nothing in
training --- a distinction worth stating because the two look identical from the outside.

\textbf{(ii) A verl-side parser recovers calls, but none of them are the tool.} We implemented a
\texttt{qwen2\_5\_coder} parser inside verl's registry accepting \texttt{<tools>} tags,
\texttt{json} fenced blocks and bare JSON. On the exact training condition (1.5B, mandatory
prompt, $n{=}100$) it lifts recoverable calls from \textbf{0/100 (hermes) to 52/100}. Of those 52,
\textbf{zero name} \texttt{run\_tests} (Table~\ref{tab:probe_categories}). Every call targets the
function the task asks the model to \emph{write} --- \texttt{can\_form\_word(\{"tiles", "word"\})},
\texttt{check\_password\_strength(\{"password"\})} --- the same signature as Llama's failure in
\S\ref{sec:llama}, now at a 5$\times$ smaller model. \textbf{Name normalisation cannot repair
this: no call carries code to normalise.}

\textbf{(iii) A role-disambiguation few-shot makes it worse, not better.} Since the failure is
semantic rather than syntactic, we targeted it directly: a system prompt containing the wrong call
and the right call side by side, with an explicit warning that the task function is not a tool.
Same model, same 100 items, same server; only the system prompt changes. The intervention made the
model \emph{more fluent at emitting the wrong call} (52 $\to$ 64 recoverable) and slightly worse at
the task (15 $\to$ 13), with the count that matters staying at zero
(Table~\ref{tab:pressure}). This is the third independent instance of the pattern in
\S\ref{sec:pressure}.

\textbf{(iv) Constrained decoding is unavailable in this stack.} Forcing the tool choice, or
schema-constrained decoding, is not exposed by verl 0.9.0's vLLM rollout path (its only
\texttt{strict} key belongs to the profiler). This remains true but was the wrong diagnosis of
what the repaired arm needed: guided decoding forces a format, whereas the successful intervention
accepts the model's existing format. See \S\ref{sec:rl}.
